\section{Results}

\subsection{Main Finding}

\textbf{Permutation equivariance is a prerequisite for learning weight-to-accuracy mappings from raw weights.}

\begin{table}[h]
\centering
\begin{tabular}{lccc}
\toprule
\textbf{N} & \textbf{Statistics $R^2$} & \textbf{NFN $R^2$} & \textbf{MLP $R^2$} \\
\midrule
100 & 0.9995 & $\sim$0.995 & $< 0$ \\
500 & 0.9995 & 0.9985 & $-1.50$ \\
5000 & 0.9995 & 0.9973 & $-1.08$ \\
\bottomrule
\end{tabular}
\end{table}

MLP fails categorically at all $N$. NFN succeeds at all $N$.

\subsection{Analysis}

The MLP faces 270K features with only thousands of samples---severely underdetermined. Each layer has $N!$ equivalent permutations. Without equivariance, the MLP cannot generalize across permutation-equivalent inputs.

NFN constrains the hypothesis space through weight sharing and architectural invariance, reducing effective dimensionality. Statistics succeeds via aggressive projection (270K $\rightarrow$ 63 features), trading expressiveness for guaranteed invariance.
